\chapter{Normalization concordance}
\label{app:normalizations}

This appendix is a firewall.  The same geometric datum is represented in the
book by a valuation, a metric weight, a model divisor, a Lagerberg potential,
a complex quasi-psh function, or a convex function.  The translations are
simple only after signs and factors have been fixed.  We collect them here in
a form that can be tested on principal divisors, the special fibre, and the
one-dimensional torus.

The entries below are conventions of this book, not claims that all the
listed theories are defined on the same objects.  A displayed comparison is
an identity only on the common locus established in the relevant comparison
chapter.  In particular, the notation \(\ddc\) belongs to complex geometry,
whereas \(\dprimeddprime\) belongs to the Chambert--Loir--Ducros calculus.
They are never identified formally.

\section{Absolute values, valuations, and scaling}
\label{sec:appJ-valuations}

\index{valuation!sign convention}
Our absolute values are multiplicative and our valuations additive:
\begin{equation*}
 v_x(f)=-\log|f(x)|.
 \tag{J.1}\label{eq:appJ-valuation-sign}
\end{equation*}
Consequently
\[
 |fg|=|f||g|,
 \qquad
 v_x(fg)=v_x(f)+v_x(g),
 \qquad
 |f+g|\leq\max\{|f|,|g|\}.
\]
A large positive value of \(v_x(f)\) means that \(f\) is small at \(x\).
For a semivaluation, \(v_x(f)=+\infty\) is equivalent to
\(f\in\Supp(x)\), or equivalently \(|f(x)|=0\).

If \(K\) is discretely valued and \(\pi\) is a uniformizer, the default
normalization is
\begin{equation*}
 v_K(\pi)=1,
 \qquad |\pi|=\rme^{-1}.
 \tag{J.2}\label{eq:appJ-uniformizer}
\end{equation*}
Replacing \(|\cdot|\) by \(|\cdot|^c\), with \(c>0\), replaces every
valuation by \(cv\).  Thus lengths in a skeleton, values of model functions,
and unnormalized log discrepancies scale by \(c\).  Intersection numbers on
an algebraic model do not scale.  A formula mixing these quantities must
therefore display the chosen normalization of the ground valuation.

Over a trivially valued field the scaling action is
\begin{equation*}
 (t\cdot v)(f)=t\,v(f),\qquad t\in\mathbf R_{>0}.
 \tag{J.3}\label{eq:appJ-scaling-action}
\end{equation*}
In multiplicative notation this is \(|f|_{t\cdot v}=|f|_v^t\).  A function
\(\phi\) on the valuation cone is homogeneous of degree one when
\(\phi(t v)=t\phi(v)\).  The trivial valuation is the cone vertex and has
value zero under every finite homogeneous function.

\begin{remark}[Two normalizations of a divisorial valuation]
If \(E\) is a prime divisor on a model, \(\ord_E\) is the integral valuation
whose value group is \(\mathbf Z\).  In a nontrivially valued model one often
uses
\[
 x_E=b_E^{-1}\ord_E,
 \qquad b_E=\ord_E(\pi),
\]
so that \(v_{x_E}(\pi)=1\).  In the trivially valued cone one normally keeps
\(\ord_E\), or fixes a slice such as \(A(v)=1\) or \(v(\mathfrak m)=1\).
These points agree only after the indicated rescaling.
\end{remark}

\section{Metric weights and tensor operations}
\label{sec:appJ-metrics}

\index{metric!additive weight}
For non-Archimedean metrics, an additive relative weight \(\phi\) is defined
by
\begin{equation*}
 \|s\|_{\phi}=\rme^{-\phi}\|s\|_0.
 \tag{J.4}\label{eq:appJ-na-weight}
\end{equation*}
Thus increasing \(\phi\) decreases the norm.  If \(\phi_i\) is a weight on
\(L_i\), the tensor-product weight on \(L_1\otimes L_2\) is
\(\phi_1+\phi_2\), and the weight on \(mL\) is \(m\phi\).  For a section
\(s\) of \(mL\), the Fubini--Study term used throughout the book is
\begin{equation*}
 m^{-1}\bigl(\log|s|+\lambda\bigr),
 \tag{J.5}\label{eq:appJ-fs-term}
\end{equation*}
with maxima of finitely many such terms.  At a valuation \(v\),
\(\log|s|(v)=-v(s)\); positive vanishing therefore lowers an FS potential.

For complex Hermitian metrics we use squared norms:
\begin{equation*}
 |s|_h^2=\rme^{-\varphi}|s|_{h_0}^2.
 \tag{J.6}\label{eq:appJ-complex-weight}
\end{equation*}
The additive weight of the underlying norm is \(\varphi/2\).  Consequently,
if a complex ray \(\varphi_t\) is compared to a non-Archimedean norm weight,
one must first check whether the source uses \(-\log|s|_h\) or
\(-\log|s|_h^2\).  In this book the complex potential is the latter.

Adding a constant has the same sign in both conventions:
\[
 \|s\|_{\phi+c}=\rme^{-c}\|s\|_\phi,
 \qquad
 |s|_{h_{\varphi+c}}^2=\rme^{-c}|s|_{h_\varphi}^2.
\]
It follows that energies satisfy \(E(\phi+c)=E(\phi)+cM\) when the total
Monge--Amp\`ere mass is \(M\).  If the energy has first been divided by
\(M\), the translation law becomes \(E(\phi+c)=E(\phi)+c\).

\section{Vertical divisors and the special-fibre test}
\label{sec:appJ-model-functions}

\index{model function!sign convention}
Let \(D\) be a vertical Cartier divisor, locally cut out by \(f=0\).  The
local generator of \(\cO(D)\) is \(f^{-1}\).  The associated model function
is
\begin{equation*}
 \phi_D=\log|f^{-1}|=-\log|f|=v_x(f).
 \tag{J.7}\label{eq:appJ-model-divisor}
\end{equation*}
Therefore
\[
 \phi_{D+D'}=\phi_D+\phi_{D'},
 \qquad
 \phi_{-D}=-\phi_D,
 \qquad
 D\geq0\Longrightarrow \phi_D\geq0.
\]
For a discretely valued field with \(v_K(\pi)=1\), the special fibre is
locally defined by \(\pi\), hence
\begin{equation*}
 \phi_{\cX_s}=1.
 \tag{J.8}\label{eq:appJ-special-fibre-test}
\end{equation*}
Equations \eqref{eq:appJ-model-divisor} and
\eqref{eq:appJ-special-fibre-test} are the two fastest sign tests in the
book.  Any formula assigning \(-1\) to the special fibre uses the opposite
metric convention.

If \(D=\operatorname{div}(g)\) is principal on a model, the associated
line bundle is algebraically trivial, but the restriction of
\(-\log|g|\) to the analytic generic fibre need not be the zero function.
One must not replace the relative intersection identities by the false
slogan that every intersection involving a principal vertical divisor
vanishes term by term.  What vanishes is the appropriate global principal
intersection after all components and horizontal contributions have been
included.

\section{Complex differential operators}
\label{sec:appJ-complex-ddc}

\index{Poincare--Lelong formula@Poincar\'e--Lelong formula!complex normalization}
The complex operator is
\begin{equation*}
 \ddc=\frac{\rmi}{2\pi}\partial\bar\partial.
 \tag{J.9}\label{eq:appJ-ddc}
\end{equation*}
With this choice, a nonzero meromorphic function satisfies
\begin{equation*}
 \ddc\log|f|^2=[\operatorname{div}(f)],
 \qquad
 \ddc\log|f|=\frac12[\operatorname{div}(f)].
 \tag{J.10}\label{eq:appJ-complex-pl}
\end{equation*}
If a local holomorphic frame \(e\) has squared norm
\(|e|_h^2=\rme^{-\varphi_e}\), then
\begin{equation*}
 c_1(L,h)=\ddc\varphi_e.
 \tag{J.11}\label{eq:appJ-complex-c1}
\end{equation*}
For a meromorphic section \(s=fe\),
\(-\log|s|_h^2=\varphi_e-\log|f|^2\), and hence
\begin{equation*}
 c_1(L,h)=\ddc\bigl(-\log|s|_h^2\bigr)
              +[\operatorname{div}(s)].
 \tag{J.12}\label{eq:appJ-complex-section-pl}
\end{equation*}
This form of Poincar\'e--Lelong is the safest one to compare with the CLD
section formula below.

If \(\theta=c_1(L,h_0)\) and \(h=h_0\rme^{-\varphi}\) in squared-norm
notation, then
\begin{equation*}
 c_1(L,h)=\theta+\ddc\varphi.
 \tag{J.13}\label{eq:appJ-complex-curvature-change}
\end{equation*}
Thus a \(\theta\)-psh function is an upper semicontinuous \(\varphi\) with
\(\theta+\ddc\varphi\geq0\).  No factor two is missing in
\eqref{eq:appJ-complex-curvature-change}; the factor occurs only when one
switches from the squared norm to the norm itself.

\section{Lagerberg and CLD conventions}
\label{sec:appJ-cld}

\index{superform!sign convention}
On \(N_{\mathbf R}\simeq\mathbf R^n\) with coordinates \(x_1,\dots,x_n\),
Lagerberg superforms use two copies of the cotangent basis.  The involution
\(J\) interchanges \(\dprime x_i\) and \(\ddprime x_i\), with the graded
signs fixed in \cref{chap:superforms}.  The positive top orientation used for integration
is
\begin{equation*}
 \tau_n=
 \dprime x_1\wedge\ddprime x_1\wedge\cdots\wedge
 \dprime x_n\wedge\ddprime x_n.
 \tag{J.14}\label{eq:appJ-lagerberg-orientation}
\end{equation*}
Equivalently, when a strongly positive \((p,p)\)-form is written with all
\(\dprime\)-factors first and all \(\ddprime\)-factors second, the positivity
generator carries
\begin{equation*}
 c_p=(-1)^{p(p-1)/2}.
 \tag{J.15}\label{eq:appJ-cp}
\end{equation*}
Strong positivity itself does not insert an additional \(c_p\) in each
rank-one factor.  This distinction prevents a double sign in even bidegree.

For a smooth function \(u\), the local real Hessian convention is
\begin{equation*}
 \dprimeddprime u
 =\sum_{i,j}\frac{\partial^2u}{\partial x_i\partial x_j}
       \dprime x_i\wedge\ddprime x_j.
 \tag{J.16}\label{eq:appJ-real-hessian}
\end{equation*}
In particular, convex functions produce positive \((1,1)\)-currents.
This is the real sign checkpoint.

\index{Poincare--Lelong formula@Poincar\'e--Lelong formula!CLD normalization}
For a nonzero meromorphic function on a boundaryless pure-dimensional
Berkovich analytic space, the CLD normalization is
\begin{equation*}
 \dprimeddprime[\log|f|]=\delta_{\operatorname{div}(f)}.
 \tag{J.17}\label{eq:appJ-cld-pl}
\end{equation*}
Since \(v(f)=-\log|f|\), it follows that
\begin{equation*}
 \dprimeddprime[v(f)]=-\delta_{\operatorname{div}(f)}.
 \tag{J.18}\label{eq:appJ-valuation-pl}
\end{equation*}
There is no contradiction with the positivity of an effective model divisor:
the current of the model metric includes its curvature datum and is not the
bare current of the principal function on an arbitrary chart.

If \(e\) is a nonvanishing local frame of a metrized line bundle, write
\(\lambda_e=-\log\|e\|\).  Then
\begin{equation*}
 c_1(L,\|\cdot\|)=\dprimeddprime\lambda_e.
 \tag{J.19}\label{eq:appJ-cld-c1}
\end{equation*}
For a meromorphic section \(s=fe\),
\(-\log\|s\|=\lambda_e-\log|f|\), and
\begin{equation*}
 c_1(L,\|\cdot\|)
 =\dprimeddprime[-\log\|s\|]
  +\delta_{\operatorname{div}(s)}.
 \tag{J.20}\label{eq:appJ-cld-section-pl}
\end{equation*}
Thus \eqref{eq:appJ-complex-section-pl} and
\eqref{eq:appJ-cld-section-pl} have exactly the same formal shape once the
complex squared norm and the non-Archimedean norm have been distinguished.

\section{Currents, cycles, and boundary signs}
\label{sec:appJ-currents}

Currents act on compactly supported test forms by the conventions of
\crefrange{chap:cld-forms}{chap:cld-bedford-taylor}.  The differential of a
current is defined by duality with
the sign that makes Stokes' formula hold.  Rather than memorize that sign in
isolation, use the following three checks:
\begin{enumerate}
\item the current associated to a smooth form differentiates to the current
  associated to its differential;
\item integration over a polyhedron satisfies
  \(\langle \dprime[\sigma],\alpha\rangle
    =\langle[\partial\sigma],\alpha\rangle\) with the induced calibration;
\item balanced internal facets cancel in a tropical cycle.
\end{enumerate}
The outward primitive normal and the induced facet weight are therefore part
of the datum.  Reversing the orientation of a cell reverses its integration
current, but it does not change a positive tropical multiplicity.

For a codimension-\(p\) cycle \(Z\), \(\delta_Z\) has bidegree \((p,p)\).
Wedge products with \(\delta_Z\) are asserted only when the pullback or
proper-intersection hypotheses of the relevant theorem hold.  The notation
\(\alpha\wedge\delta_Z\) is not a license to multiply arbitrary singular
currents.

\section{Model intersections and Chambert--Loir measures}
\label{sec:appJ-model-ma}

\index{Chambert--Loir measure!normalization}
Let \(X\) be an \(n\)-dimensional proper variety over a nontrivially valued
field, let \(\cX\) be a normal model, and let
\(\cL_1,\dots,\cL_n\) be model line bundles.  Write
\(\cX_s=\sum_E b_EE\) and let \(x_E=b_E^{-1}\ord_E\) be the corresponding
normalized divisorial point.  The unnormalized mixed Chambert--Loir measure is
\begin{equation*}
 c_1(\cL_1)\cdots c_1(\cL_n)
 =\sum_E b_E
  \bigl(\cL_1|_E\cdots\cL_n|_E\bigr)\delta_{x_E}.
 \tag{J.21}\label{eq:appJ-cl-measure}
\end{equation*}
Its total mass is the generic-fibre intersection number
\begin{equation*}
 \int_{X^{\An}}c_1(\cL_1)\cdots c_1(\cL_n)
 =(L_1\cdots L_n)_X.
 \tag{J.22}\label{eq:appJ-cl-total-mass}
\end{equation*}
The multiplicity \(b_E\) and the normalization of \(x_E\) occur together.
Using \(\ord_E\) in place of \(x_E\) without modifying the coefficient moves
the measure to a differently scaled valuation cone.

For one model line bundle \(\cL\), the Monge--Amp\`ere notation is
\begin{equation*}
 \MA(\phi_{\cL})=c_1(\cL)^n,
 \qquad
 \MA(\phi_{\cL})(X^{\An})=(L^n).
 \tag{J.23}\label{eq:appJ-model-ma}
\end{equation*}
When \((L^n)>0\), the probability version is
\begin{equation*}
 \MA_{\mathrm{prob}}(\phi)
  =(L^n)^{-1}\MA(\phi).
 \tag{J.24}\label{eq:appJ-prob-ma}
\end{equation*}
The probability notation is never used in a zero-volume class.

\section{BFJ and BJ energy normalizations}
\label{sec:appJ-energy}

Fix a reference psh metric \(\phi_0\) and let
\(V=(L^n)>0\).  The unnormalized mixed Aubin--Mabuchi energy is
\begin{equation*}
 E_{\phi_0}(\phi)
 =\frac1{n+1}\sum_{j=0}^n
   \int_{X^{\An}}(\phi-\phi_0)\,
   \MA(\phi^{(j)},\phi_0^{(n-j)}).
 \tag{J.25}\label{eq:appJ-energy}
\end{equation*}
It obeys
\begin{equation*}
 E_{\phi_0}(\phi_0)=0,
 \qquad
 E_{\phi_0}(\phi+c)=E_{\phi_0}(\phi)+cV.
 \tag{J.26}\label{eq:appJ-energy-translation}
\end{equation*}
Some sources divide \eqref{eq:appJ-energy} by \(V\).  Their energy translates
by \(c\), and their first variation is integration against the probability
Monge--Amp\`ere measure.  Every comparison theorem in the main text states
which version is used.

The cocycle is
\begin{equation*}
 E_{\phi_0}(\phi)-E_{\phi_0}(\psi)
 =\frac1{n+1}\sum_{j=0}^n
   \int(\phi-\psi)\,
   \MA(\phi^{(j)},\psi^{(n-j)}).
 \tag{J.27}\label{eq:appJ-energy-cocycle}
\end{equation*}
In particular, \(\phi\leq\psi\) implies \(E(\phi)\leq E(\psi)\).  Along a
differentiable path,
\begin{equation*}
 \frac{\rd}{\rd t}E(\phi_t)
 =\int\dot\phi_t\,\MA(\phi_t).
 \tag{J.28}\label{eq:appJ-energy-variation}
\end{equation*}
These two tests fix the sign of energy.

For a probability measure \(\mu\), the Legendre energy used in the BJ part is
schematically
\begin{equation*}
 E^\vee(\mu)=
 \sup_{\phi}\bigl(E_{\mathrm{prob}}(\phi)-\textstyle\int\phi\,\rd \mu\bigr).
 \tag{J.29}\label{eq:appJ-dual-energy}
\end{equation*}
It is translation invariant because both terms change by the same constant.
With unnormalized measures of mass \(V\), both terms must instead be
unnormalized.  Mixing a probability \(\mu\) with the unnormalized energy
breaks translation invariance and is therefore immediately detectable.

\section{Complex nonpluripolar mass}
\label{sec:appJ-complex-mass}

Let \(X\) be a compact K\"ahler space of dimension \(n\), let \(\theta\) be
a smooth closed real \((1,1)\)-form, and let \(\varphi\in\PSH(X,\theta)\).
The nonpluripolar product is written
\begin{equation*}
 \langle(\theta+\ddc\varphi)^n\rangle.
 \tag{J.30}\label{eq:appJ-npp}
\end{equation*}
Its mass can be strictly smaller than \(\vol(\theta)\).  Full Monge--Amp\`ere
mass means equality with \(\vol(\theta)\), not probability normalization.
When \(\vol(\theta)=0\), division by the volume is forbidden; the inverse
systems in \cref{part:complex} are formulated without such a division.

If \(\varphi\) and \(\psi\) have the same singularity type, their relative
energy is invariant under adding the same constant to both potentials.  A
test curve or ray may use an energy normalized by \(\vol(\theta)\), but only
on the positive-volume locus.  This is why the zero-volume theory in
\crefrange{chap:transcendental-metrics}{chap:transcendental-okounkov-dh} is
kept separate.

\section{Toric coordinates and real Monge--Ampere}
\label{sec:appJ-toric}

\index{toric geometry!sign convention}
For a character \(\chi^m\) on an algebraic torus, our tropical coordinate is
\begin{equation*}
 u(x)=-\log|\chi(x)|,
 \qquad
 \langle m,u(x)\rangle=-\log|\chi^m(x)|.
 \tag{J.31}\label{eq:appJ-toric-coordinate}
\end{equation*}
If a toric metric is represented by its additive weight
\(-\log\|s\|\), then the recession function attached to a polytope \(P\) is
\begin{equation*}
 \check h_P(u)=h_P(-u)=\sup_{m\in P}\langle m,-u\rangle.
 \tag{J.32}\label{eq:appJ-toric-recession}
\end{equation*}
Accordingly the subgradient image lies in \(-P\), not \(P\).  Authors using
\(\log|\chi|\) as tropical coordinate obtain the opposite sign.

For a convex function \(g\) on \(N_{\mathbf R}\), the Alexandrov measure is
normalized by lattice Lebesgue measure:
\begin{equation*}
 \MA_{\mathbf R}(g)(B)
 =\operatorname{Leb}_M(\partial g(B)).
 \tag{J.33}\label{eq:appJ-alexandrov}
\end{equation*}
For a smooth strictly convex \(g\), this becomes
\(\det(D^2g)\,\rd u\).  Algebraic intersection numbers introduce the factor
\(n!\): for a lattice polytope \(P\),
\begin{equation*}
 (L_P^n)=n!\operatorname{vol}_M(P).
 \tag{J.34}\label{eq:appJ-toric-degree}
\end{equation*}
Thus a toric comparison may use \(\MA_{\mathbf R}\),
\(n!\MA_{\mathbf R}\), or a probability measure.  The total-mass test
\eqref{eq:appJ-toric-degree} identifies the choice.

\section{Rays, test configurations, and Duistermaat--Heckman measures}
\label{sec:appJ-rays}

Let \(\varphi_t\) be a complex weak geodesic ray.  The book parametrizes it
for \(t\geq0\); translating the associated non-Archimedean metric by \(c\)
adds the linear function \(ct\) to the ray.  Consequently the asymptotic
energy slope changes by \(cV\) in the unnormalized convention.

A decreasing test curve \(\tau\mapsto\psi_\tau\) and its ray are related by
the Legendre formulas
\begin{equation*}
 \varphi_t=\uscsup[\tau\in\mathbf R]
                 (\psi_\tau+t\tau),
 \qquad
 \psi_\tau=\inf_{t>0}(\varphi_t-t\tau),
 \tag{J.35}\label{eq:appJ-legendre-rays}
\end{equation*}
with the separate \(t=0\) convention stated in \cref{chap:rays-test-curves}.  Reversing
\(t\tau\) to \(-t\tau\) reflects the Duistermaat--Heckman measure.  Moment
identities therefore use the sign in \eqref{eq:appJ-legendre-rays}.

The unnormalized Duistermaat--Heckman measure has total mass \(V\).  Its
probability normalization is \(V^{-1}\DHm\).  If \(\lambda\) denotes the
filtration weight, then translating all weights by \(c\) pushes the measure
forward by \(\lambda\mapsto\lambda+c\).  This translation test catches both
sign and normalization mistakes in slope formulae.

For a normal test configuration \((\cX,\cL)\) over \(\mathbf A^1\), the
compactification over \(\mathbf P^1\) is used for intersection numbers.  The
choice \(\phi_{\cX_0}=1\) agrees with
\eqref{eq:appJ-special-fibre-test}; a source using \(-\log|\tau|^{-1}\) as
parameter may instead attach \(-1\) to the central fibre.  We translate such
formulae before importing them.

\section{Log discrepancy, entropy, and stability functionals}
\label{sec:appJ-stability}

On a trivially valued cone, log discrepancy is homogeneous:
\begin{equation*}
 A_X(tv)=tA_X(v).
 \tag{J.36}\label{eq:appJ-discrepancy-homogeneous}
\end{equation*}
For a divisorial valuation \(v=c\ord_E\) over a normal \(\mathbf Q\)-Gorenstein
variety,
\begin{equation*}
 A_X(v)=c\bigl(1+\ord_E(K_Y-\pi^*K_X)\bigr).
 \tag{J.37}\label{eq:appJ-div-discrepancy}
\end{equation*}
For a pair \((X,B)\), replace the relative canonical divisor by
\(K_Y-\pi^*(K_X+B)\).  The coefficient \(1\) is part of log discrepancy;
omitting it gives the discrepancy coefficient instead.

Entropy of a probability measure is
\begin{equation*}
 \Ent_X(\mu)=\int A_X\,\rd \mu\in[0,+\infty].
 \tag{J.38}\label{eq:appJ-entropy}
\end{equation*}
For measures of mass \(V\), the unnormalized entropy is the same integral and
has mass scale \(V\).  Ding and Mabuchi functionals in \cref{chap:bj-stability} use a
declared normalization of energy and entropy together.  Their coercivity
constants are meaningless unless the same normalization is used on both
sides of the inequality.

The complex-ray formalism uses a different slice.  On valuations of
\(X\times\mathbf C\) normalized by \(v(z)=1\), the total discrepancy
\(A_{X\times\mathbf C}\) takes the value \(1\) at the product valuation.
The function compatible with entropy slopes is therefore the reduced
discrepancy \(A_{X\times\mathbf C}-1\), which vanishes there.  It must not be
identified with either total discrepancy on that slice or the homogeneous
cone function \(A_X\) in \eqref{eq:appJ-discrepancy-homogeneous}; the
constant-ray test distinguishes the three conventions.

The Donaldson--Futaki invariant depends on the sign chosen for the
\(\mathbf G_m\)-weight.  Our convention is tied to
\eqref{eq:appJ-fs-term} and \eqref{eq:appJ-legendre-rays}; the weight
distribution is the Duistermaat--Heckman measure appearing there.  When a
source acts by \(\tau\mapsto\tau^{-1}\), all weights and the first moment are
reflected before comparison.

\section{A compact translation table}
\label{sec:appJ-table}

The following table is a mnemonic, not a substitute for the domain
hypotheses in the main text.
\begin{center}
\small
\begin{tabular}{@{}p{.22\textwidth}p{.31\textwidth}p{.34\textwidth}@{}}
\toprule
Datum & Book convention & Diagnostic test \\
\midrule
Absolute value and valuation
 & \(v=-\log|\cdot|\)
 & \(|\pi|=\rme^{-1}\Rightarrow v(\pi)=1\) \\
NA metric weight
 & \(\|s\|_\phi=\rme^{-\phi}\|s\|_0\)
 & larger \(\phi\) means smaller norm \\
Complex metric weight
 & \(|s|_h^2=\rme^{-\varphi}|s|_{h_0}^2\)
 & norm weight is \(\varphi/2\) \\
Vertical divisor
 & \(\phi_D=v(f_D)\)
 & \(\phi_{\cX_s}=1\) \\
Complex Poincar\'e--Lelong
 & \(\ddc\log|f|^2=[\operatorname{div}(f)]\)
 & \(\ddc=(\rmi/2\pi)\partial\bar\partial\) \\
CLD Poincar\'e--Lelong
 & \(\dprimeddprime\log|f|=\delta_{\operatorname{div}(f)}\)
 & \(\dprimeddprime v(f)=-\delta_{\operatorname{div}(f)}\) \\
Top Lagerberg orientation
 & \(\tau_n=\bigwedge_i(\dprime x_i\wedge\ddprime x_i)\)
 & convex Hessian is positive \\
Monge--Amp\`ere mass
 & \(\MA(\phi)(X)=(L^n)\)
 & probability version divides by \((L^n)\) \\
Toric coordinate
 & \(u=-\log|\chi|\)
 & recession \(h_P(-u)\), image \(-P\) \\
Complex--NA Legendre sign
 & \(\varphi_t=\sup_\tau^*(\psi_\tau+t\tau)\)
 & weight translation gives \(ct\) \\
Log discrepancy
 & \(A(c\ord_E)=c(1+\ord_EK_{Y/X})\)
 & homogeneous of degree one \\
\bottomrule
\end{tabular}
\end{center}

\section{The normalization unit tests}
\label{sec:appJ-unit-tests}

Before accepting a cross-theory formula, apply the following tests in order.
They are deliberately redundant.

\begin{enumerate}
\item \emph{Constant test.}  Replace a potential by \(\phi+c\).  Curvature
  and Monge--Amp\`ere measure must remain unchanged.  An unnormalized energy
  changes by \(cV\), a normalized energy by \(c\), and a translation-invariant
  dual functional does not change.

\item \emph{Principal-divisor test.}  Substitute a meromorphic function and
  use \eqref{eq:appJ-complex-pl} or \eqref{eq:appJ-cld-pl}.  The formula must
  distinguish \(\log|f|\) from \(-\log|f|=v(f)\).

\item \emph{Special-fibre test.}  Put \(D=\cX_s\) and
  \(v_K(\pi)=1\).  The associated model function must be \(+1\).

\item \emph{Total-mass test.}  Integrate the proposed Monge--Amp\`ere
  identity.  Both sides must equal \((L^n)\), \(\vol(\theta)\), or \(1\),
  according to the declared convention.

\item \emph{Toric interval test.}  On \(\mathbf P^1\), use a lattice segment
  of length \(d\).  The algebraic degree is \(d\); the subgradient image is
  the reflected interval dictated by \eqref{eq:appJ-toric-recession}.

\item \emph{Scaling test.}  Replace \(v\) by \(tv\).  Homogeneous
  potentials and log discrepancies scale by \(t\), while projective
  intersection numbers do not.

\item \emph{Zero-volume test.}  Put \(\vol(\theta)=0\).  Any formula that
  divides by the volume or invokes a probability Monge--Amp\`ere measure has
  left its domain of definition.

\item \emph{Support test.}  Evaluate at a semivaluation with proper support.
  A function vanishing on that support has value \(+\infty\), and its
  logarithmic norm is \(-\infty\).  Formulae written only on the
  function-field valuation locus must say so.
\end{enumerate}

Passing all eight tests does not prove a comparison theorem, but failing one
of them disproves the proposed normalization.  They will be cited repeatedly
in \cref{part:comparisons}.

\section*{Dependency and source note}

The valuation and metric conventions are fixed in
\cref{chap:four-entrances,chap:valuations,chap:metrics,chap:trivial-analytification}.
The Lagerberg orientation and positivity signs are proved in
\crefrange{chap:superforms}{chap:tropical-integration};
the CLD Poincar\'e--Lelong and Chern-current formulae are established in
\cref{chap:cld-poincare-lelong,chap:cld-metrized-bundles}.  Model intersection
measures are treated in
\cref{chap:metrics,chap:cld-metrized-bundles,chap:bfj-ma,chap:bj-fs-pl-psh}.
The BFJ energy conventions are fixed in
\cref{chap:bfj-energy,chap:bfj-finite-energy-measures}, and the BJ conventions
in \crefrange{chap:bj-energy}{chap:bj-strong-topologies}.  Complex nonpluripolar
products, test curves, and transcendental inverse systems are normalized in
\crefrange{chap:complex-prerequisites}{chap:transcendental-okounkov-dh}.  The
comparison statements themselves belong to
\crefrange{chap:master-dictionary}{chap:ytd-bridge}; this appendix
only makes their bookkeeping explicit.
